% Auto-generated complete chapter from 29_portfolio_case_study_benchmark.md
\chapter{29\_portfolio\_case\_study\_benchmark}
\label{complete:root_037}

\begin{sourcemeta}
\textbf{Fuente:} \href{file:///E:/Laboral/29_portfolio_case_study_benchmark.md}{29\_portfolio\_case\_study\_benchmark.md}\\
\textbf{Seccion:} root\\
\textbf{Estado:} canonical\\
\textbf{Rol:} Modulo principal\\
\textbf{Lineas originales:} 1011\\
\textbf{Ultima modificacion:} 2026-06-11 18:08\\
\textbf{Deduplicacion conservadora:} 0 bloques repetidos omitidos; 0 caracteres repetidos no reimpresos.
\end{sourcemeta}

\section*{Resumen de orientacion}
This module benchmarks public case studies relevant to \textbf{TwinSight X500} and Alexander Woodcock Salomón’s positioning as:

\section*{Contenido completo depurado}
\addcontentsline{toc}{section}{Contenido completo depurado}




\subsection{Module status}




\begin{mdcode}
Bloque de codigo: text
Module: 29_portfolio_case_study_benchmark.md
Input task: A5_technical_visualization_case_studies
Date generated: 2026-06-11
Status: Partial but usable
Source type: public web case studies
\end{mdcode}




\subsection{Scope}




This module benchmarks public case studies relevant to \textbf{TwinSight X500} and Alexander Woodcock Salomón’s positioning as:




\begin{mdcode}
Bloque de codigo: text
Real-Time 3D Developer / Unity Technical Artist
focused on Unity WebGL, CAD-to-realtime optimization, interactive technical visualization, and technical UI.
\end{mdcode}




The source collection focused on case studies involving:




\begin{itemize}
\item Unity WebGL;
\item browser-based interactive 3D;
\item technical visualization;
\item digital twins;
\item XR / VR training;
\item product configurators;
\item CAD / BIM / reality-capture workflows;
\item industrial visualization;
\item spatial onboarding;
\item training simulations.
\end{itemize}




\subsection{Important limitation}




The A5 collection produced \textbf{7 usable detailed case studies}, not the original target of 20–30.




This module is still useful because the collected examples have enough structure to extract case-study patterns for TwinSight. However, it should be treated as:




\begin{mdcode}
Bloque de codigo: text
Case-study benchmark v1
\end{mdcode}




A later expansion can add:




\begin{mdcode}
Bloque de codigo: text
29B_demo_reels_benchmark.md
29C_artstation_breakdown_benchmark.md
29D_case_study_benchmark_expansion.md
\end{mdcode}




\medskip\hrule\medskip




\section{1. Source inventory}




\subsection{1.1 Case studies collected}




\begin{center}
\scriptsize
\begin{longtable}{>{\raggedright\arraybackslash}p{0.225\linewidth} >{\raggedright\arraybackslash}p{0.225\linewidth} >{\raggedright\arraybackslash}p{0.225\linewidth} >{\raggedright\arraybackslash}p{0.225\linewidth}}
\toprule
\# & Case study & Company & Source link \\
\midrule
\endfirsthead
\toprule
\# & Case study & Company & Source link \\
\midrule
\endhead
1 & WebXR Product Configurator & HapzXR & [Open](\url{https://hapzxr.com/case-studies/webxr-product-configurator/}) \\
2 & Digital Twin of Production Facility & HapzXR & [Open](\url{https://hapzxr.com/case-studies/digital-twin-of-production-facility/}) \\
3 & BIM Site Overlay Visualization & HapzXR & [Open](\url{https://hapzxr.com/case-studies/bim-site-overlay-visualization/}) \\
4 & Construction Confined Space Simulation & HapzXR & [Open](\url{https://hapzxr.com/case-studies/construction-confined-space-simulation/}) \\
5 & Work Spatial — AI-Driven Onboarding Platform & Virtual Verse Studio & [Open](\url{https://virtualverse.studio/projects/work-spatial/}) \\
6 & NBK Virtugate — National Bank of Kuwait & Virtual Verse Studio & [Open](\url{https://virtualverse.studio/projects/nbk-virtugate/}) \\
7 & Empathy Lab — UK Rail Sector Training & Virtual Verse Studio & [Open](\url{https://virtualverse.studio/projects/empathy-lab/}) \\
\bottomrule
\end{longtable}
\end{center}




\subsection{1.2 Supporting source pages}




\begin{center}
\scriptsize
\begin{longtable}{>{\raggedright\arraybackslash}p{0.305\linewidth} >{\raggedright\arraybackslash}p{0.305\linewidth} >{\raggedright\arraybackslash}p{0.305\linewidth}}
\toprule
Source & Use & Link \\
\midrule
\endfirsthead
\toprule
Source & Use & Link \\
\midrule
\endhead
Unity Industry & Context for Unity in industry & [Open](\url{https://unity.com/industry}) \\
Unity Digital Twin Applications & Digital twin context & [Open](\url{https://unity.com/topics/digital-twin-applications-and-use-cases}) \\
Virtual Verse Unity WebGL page & WebGL service positioning & [Open](\url{https://virtualverse.studio/unity-webgl/}) \\
Prevu3D Case Studies & Digital twin expansion source & [Open](\url{https://www.prevu3d.com/digital-twin-case-studies/}) \\
\bottomrule
\end{longtable}
\end{center}




\medskip\hrule\medskip




\section{2. Benchmark table — identity and relevance}




\begin{center}
\scriptsize
\begin{longtable}{>{\raggedright\arraybackslash}p{0.225\linewidth} >{\raggedright\arraybackslash}p{0.225\linewidth} >{\raggedright\arraybackslash}p{0.225\linewidth} >{\raggedright\arraybackslash}p{0.225\linewidth}}
\toprule
Case & Sector & Project type & Relevance \\
\midrule
\endfirsthead
\toprule
Case & Sector & Project type & Relevance \\
\midrule
\endhead
WebXR Product Configurator & Industrial equipment & Web 3D configurator & Very high \\
Digital Twin Facility & Industrial operations & Real-time digital twin & Very high \\
BIM Site Overlay & Construction & BIM / VR overlay & High \\
Confined Space Simulation & Construction safety & VR procedure training & Medium-high \\
Work Spatial & Enterprise onboarding & WebGL XR platform & High \\
NBK Virtugate & Banking onboarding & WebGL + VR learning & High \\
Empathy Lab & Rail training & VR soft-skill simulation & Medium \\
\bottomrule
\end{longtable}
\end{center}




\subsection{Relevance definition}




\begin{center}
\scriptsize
\begin{longtable}{>{\raggedright\arraybackslash}p{0.455\linewidth} >{\raggedright\arraybackslash}p{0.455\linewidth}}
\toprule
Level & Meaning \\
\midrule
\endfirsthead
\toprule
Level & Meaning \\
\midrule
\endhead
Very high & Directly maps to TwinSight: WebGL, technical visualization, digital twin, product inspection, or browser-based 3D \\
High & Strongly useful for case-study structure or interaction design \\
Medium-high & Useful for training/procedural logic \\
Medium & Useful for presentation pattern, but weaker domain match \\
\bottomrule
\end{longtable}
\end{center}




\medskip\hrule\medskip




\section{3. Detailed case-study notes}




\subsection{3.1 WebXR Product Configurator — HapzXR}




\textbf{Source:} [\url{https://hapzxr.com/case-studies/webxr-product-configurator/}](\url{https://hapzxr.com/case-studies/webxr-product-configurator/})




\subsubsection{Confirmed facts}




\begin{center}
\scriptsize
\begin{longtable}{>{\raggedright\arraybackslash}p{0.455\linewidth} >{\raggedright\arraybackslash}p{0.455\linewidth}}
\toprule
Field & Data \\
\midrule
\endfirsthead
\toprule
Field & Data \\
\midrule
\endhead
Source type & Company case study \\
Company & HapzXR \\
Region & India / enterprise XR services \\
Industry & Industrial equipment \\
Technology & WebGL, Three.js, WebAR \\
Project type & Browser-based 3D product configurator \\
Access model & Browser + WebAR \\
Metrics & 3× engagement; 45\% more qualified leads; 2-week shorter pre-order cycle \\
Confidence & High \\
\bottomrule
\end{longtable}
\end{center}




\subsubsection{Case structure}




\begin{mdcode}
Bloque de codigo: text
Headline
Short value proposition
Challenge
Solution
Results
Project overview
Snapshot
Delivery flow
Related work
Call to action
\end{mdcode}




\subsubsection{Problem pattern}




The buyer could not evaluate industrial equipment physically. Static catalogs and 2D renderings failed to communicate scale, detail, and configuration.




\subsubsection{Solution pattern}




A browser-based 3D configurator allowed users to:




\begin{itemize}
\item rotate;
\item zoom;
\item customize finishes;
\item inspect internal components;
\item use WebAR placement;
\item send configuration data to CRM.
\end{itemize}




\subsubsection{Direct relevance to TwinSight}




This is the closest collected analogue to TwinSight’s \textbf{browser-based inspection and technical product understanding}. It shows how a 3D viewer becomes stronger when framed as a decision-support or sales-support tool, not just a model viewer.




\subsubsection{What Alexander can adapt}




\begin{itemize}
\item Frame TwinSight as an \textbf{interactive technical understanding tool}.
\item Use “2D support was insufficient” as the problem.
\item Emphasize browser access.
\item Show interaction modes clearly.
\item Include measurable outcomes if defensible.
\item Present features as user-facing workflows, not just implementation details.
\end{itemize}




\subsubsection{What not to copy}




\begin{itemize}
\item Do not claim CRM integration.
\item Do not claim WebAR unless implemented.
\item Do not claim business conversion metrics.
\item Do not overstate industrial client impact.
\end{itemize}




\medskip\hrule\medskip




\subsection{3.2 Digital Twin of Production Facility — HapzXR}




\textbf{Source:} [\url{https://hapzxr.com/case-studies/digital-twin-of-production-facility/}](\url{https://hapzxr.com/case-studies/digital-twin-of-production-facility/})




\subsubsection{Confirmed facts}




\begin{center}
\scriptsize
\begin{longtable}{>{\raggedright\arraybackslash}p{0.455\linewidth} >{\raggedright\arraybackslash}p{0.455\linewidth}}
\toprule
Field & Data \\
\midrule
\endfirsthead
\toprule
Field & Data \\
\midrule
\endhead
Source type & Company case study \\
Company & HapzXR \\
Industry & Chemical processing / facility operations \\
Technology & Unity / Unreal listed; IoT; web; VR \\
Project type & Real-time digital twin \\
Access model & Web dashboard + VR \\
Metrics & 30\% downtime reduction; 50\% planning-time reduction \\
Confidence & High for page content; medium for stack specificity \\
\bottomrule
\end{longtable}
\end{center}




\subsubsection{Case structure}




\begin{mdcode}
Bloque de codigo: text
Title
Value proposition
Challenge
Solution
Results
Project overview
Snapshot
Outcome summary
Delivery flow
CTA
\end{mdcode}




\subsubsection{Problem pattern}




Monitoring systems lacked spatial context. Operators could see alerts but not easily understand where those alerts belonged in the facility.




\subsubsection{Solution pattern}




A navigable 3D facility twin connected equipment to:




\begin{itemize}
\item status;
\item maintenance history;
\item predictive health indicators;
\item IoT sensor feeds;
\item web monitoring;
\item VR planning.
\end{itemize}




\subsubsection{Direct relevance to TwinSight}




TwinSight is not a live digital twin, but it is \textbf{digital-twin-adjacent} because it maps a real physical product into an interactive 3D structure for understanding spatial relationships.




\subsubsection{What Alexander can adapt}




\begin{itemize}
\item Use the term \textbf{digital-twin-adjacent technical visualization} carefully.
\item Show how each component maps to real-world parts.
\item Add a “future extension” section for sensor data, maintenance status, or assembly state.
\item Emphasize spatial context as the core value.
\item Structure the case study around “2D documentation lacks spatial context.”
\end{itemize}




\subsubsection{What not to copy}




\begin{itemize}
\item Do not claim live IoT.
\item Do not claim predictive maintenance.
\item Do not claim production facility scale.
\item Do not call TwinSight a full digital twin unless live state/data integration is added.
\end{itemize}




\medskip\hrule\medskip




\subsection{3.3 BIM Site Overlay Visualization — HapzXR}




\textbf{Source:} [\url{https://hapzxr.com/case-studies/bim-site-overlay-visualization/}](\url{https://hapzxr.com/case-studies/bim-site-overlay-visualization/})




\subsubsection{Confirmed facts}




\begin{center}
\scriptsize
\begin{longtable}{>{\raggedright\arraybackslash}p{0.455\linewidth} >{\raggedright\arraybackslash}p{0.455\linewidth}}
\toprule
Field & Data \\
\midrule
\endfirsthead
\toprule
Field & Data \\
\midrule
\endhead
Source type & Company case study \\
Company & HapzXR \\
Industry & Construction / AEC \\
Technology & Unity, BIM data layering, VR review \\
Project type & BIM / site overlay visualization \\
Access model & VR review \\
Metrics & Faster decisions; earlier mismatch detection \\
Confidence & Medium-high \\
\bottomrule
\end{longtable}
\end{center}




\subsubsection{Problem pattern}




Stakeholders lacked a unified way to compare site reality against BIM plans.




\subsubsection{Solution pattern}




Immersive BIM walkthroughs overlaid site imagery and video onto BIM models to compare planned and actual states.




\subsubsection{Direct relevance to TwinSight}




The equivalent idea for TwinSight would be overlaying or connecting:




\begin{itemize}
\item manufacturer diagrams;
\item assembly instructions;
\item labels;
\item inspection zones;
\item exploded-view states;
\item cross-section states;
\item real part references.
\end{itemize}




\subsubsection{What Alexander can adapt}




\begin{itemize}
\item Add a case-study section called \textbf{From 2D support to spatial comparison}.
\item Use before/after visuals:
\item original 2D/manual support;
\item optimized 3D viewer;
\item exploded view;
\item clipping;
\item labels.
\item Present each technical mode as a “comparison aid.”
\end{itemize}




\subsubsection{What not to copy}




\begin{itemize}
\item Do not claim BIM.
\item Do not claim site photo alignment.
\item Do not claim construction domain experience unless applying broadly to visualization.
\end{itemize}




\medskip\hrule\medskip




\subsection{3.4 Construction Confined Space Simulation — HapzXR}




\textbf{Source:} [\url{https://hapzxr.com/case-studies/construction-confined-space-simulation/}](\url{https://hapzxr.com/case-studies/construction-confined-space-simulation/})




\subsubsection{Confirmed facts}




\begin{center}
\scriptsize
\begin{longtable}{>{\raggedright\arraybackslash}p{0.455\linewidth} >{\raggedright\arraybackslash}p{0.455\linewidth}}
\toprule
Field & Data \\
\midrule
\endfirsthead
\toprule
Field & Data \\
\midrule
\endhead
Source type & Company case study \\
Company & HapzXR \\
Industry & Construction safety \\
Technology & Unity, Quest/PCVR, instructor review \\
Project type & VR procedural training \\
Access model & VR \\
Metrics & Procedural consistency and confidence; no hard numeric metric \\
Confidence & Medium-high \\
\bottomrule
\end{longtable}
\end{center}




\subsubsection{Problem pattern}




Classroom briefings did not create enough readiness for high-risk confined-space operations.




\subsubsection{Solution pattern}




A simulation-first training flow guided users through:




\begin{itemize}
\item permit verification;
\item pre-entry checks;
\item hazard controls;
\item communication;
\item emergency escalation;
\item guided and assessment modes.
\end{itemize}




\subsubsection{Direct relevance to TwinSight}




TwinSight is not a safety-training system, but it has strong overlap with procedural understanding and guided technical comprehension.




\subsubsection{What Alexander can adapt}




\begin{itemize}
\item Add a \textbf{task-based walkthrough} to the case study:
\item identify component;
\item inspect assembly relation;
\item activate exploded view;
\item use clipping;
\item compare with 2D support.
\item Show how the system supports procedural reasoning.
\item Connect SUS/NASA-TLX evaluation to training-system validation language.
\end{itemize}




\subsubsection{What not to copy}




\begin{itemize}
\item Do not claim VR safety training.
\item Do not claim procedural certification.
\item Do not use “safety outcomes” unless evaluated.
\end{itemize}




\medskip\hrule\medskip




\subsection{3.5 Work Spatial — AI-Driven Onboarding Platform — Virtual Verse Studio}




\textbf{Source:} [\url{https://virtualverse.studio/projects/work-spatial/}](\url{https://virtualverse.studio/projects/work-spatial/})




\subsubsection{Confirmed facts}




\begin{center}
\scriptsize
\begin{longtable}{>{\raggedright\arraybackslash}p{0.455\linewidth} >{\raggedright\arraybackslash}p{0.455\linewidth}}
\toprule
Field & Data \\
\midrule
\endfirsthead
\toprule
Field & Data \\
\midrule
\endhead
Source type & Portfolio case study \\
Company & Virtual Verse Studio \\
Region & Egypt / USA \\
Industry & Enterprise onboarding \\
Technology & WebGL, GenAI, Web3 \\
Project type & Browser-based XR onboarding MVP \\
Access model & WebGL \\
Metrics & On-time MVP; investor/client demos \\
Confidence & High \\
\bottomrule
\end{longtable}
\end{center}




\subsubsection{Case structure}




\begin{mdcode}
Bloque de codigo: text
Breadcrumb
Tags
Title
One-sentence value proposition
Image
Overview
Challenge
Solution
Key features
Results & impact
Client quote
Links
More portfolio items
CTA
\end{mdcode}




\subsubsection{Problem pattern}




A startup needed a fast MVP combining multiple emerging technologies.




\subsubsection{Solution pattern}




The studio delivered a browser-based 3D workspace with:




\begin{itemize}
\item avatar navigation;
\item generative AI;
\item Web3 identity;
\item modular architecture.
\end{itemize}




\subsubsection{Direct relevance to TwinSight}




This case is relevant for \textbf{portfolio structure} more than domain. It shows a concise, recruiter-readable format.




\subsubsection{What Alexander can adapt}




\begin{itemize}
\item Use short tags at the top:
\item \texttt{Unity WebGL}
\item \texttt{Technical Visualization}
\item \texttt{CAD-to-Realtime}
\item \texttt{Drone Assembly}
\item Use a one-sentence project value proposition.
\item Add a \textbf{Key Features} section before deep technical details.
\item Add a \textbf{Results \& Validation} section using SUS/NASA-TLX.
\item Add direct links to demo, GitHub, video, and thesis if public.
\end{itemize}




\subsubsection{What not to copy}




\begin{itemize}
\item Do not include AI/Web3 as primary unless relevant.
\item Do not turn TwinSight into a hype-stack case study.
\item Do not overfocus on MVP speed if the project was academic and thesis-driven.
\end{itemize}




\medskip\hrule\medskip




\subsection{3.6 NBK Virtugate — National Bank of Kuwait — Virtual Verse Studio}




\textbf{Source:} [\url{https://virtualverse.studio/projects/nbk-virtugate/}](\url{https://virtualverse.studio/projects/nbk-virtugate/})




\subsubsection{Confirmed facts}




\begin{center}
\scriptsize
\begin{longtable}{>{\raggedright\arraybackslash}p{0.455\linewidth} >{\raggedright\arraybackslash}p{0.455\linewidth}}
\toprule
Field & Data \\
\midrule
\endfirsthead
\toprule
Field & Data \\
\midrule
\endhead
Source type & Portfolio case study \\
Company & Virtual Verse Studio \\
Region & Egypt / USA \\
Industry & Banking / corporate onboarding \\
Technology & Unity WebGL + VR \\
Project type & Gamified onboarding \\
Access model & Browser + VR \\
Metrics & Improved comprehension and engagement; no numeric data \\
Confidence & High \\
\bottomrule
\end{longtable}
\end{center}




\subsubsection{Problem pattern}




Traditional employee orientation was passive, resource-heavy, and inconsistent across locations.




\subsubsection{Solution pattern}




A custom WebGL environment allowed employees to explore departments, history, tools, and branded spaces through avatars.




\subsubsection{Direct relevance to TwinSight}




This example is relevant because TwinSight also converts passive/2D information into an active interactive learning environment.




\subsubsection{What Alexander can adapt}




\begin{itemize}
\item Use the contrast:
\item passive documentation;
\item active spatial exploration.
\item Emphasize browser-first deployment.
\item Present technical information as an interactive environment.
\item Use “no install required” if true for the deployed build.
\item Include “desktop browser + future XR optional” only if framed as future.
\end{itemize}




\subsubsection{What not to copy}




\begin{itemize}
\item Do not over-gamify TwinSight unless game mechanics exist.
\item Do not claim enterprise deployment.
\item Do not claim engagement improvement unless measured.
\end{itemize}




\medskip\hrule\medskip




\subsection{3.7 Empathy Lab — UK Rail Sector Training — Virtual Verse Studio}




\textbf{Source:} [\url{https://virtualverse.studio/projects/empathy-lab/}](\url{https://virtualverse.studio/projects/empathy-lab/})




\subsubsection{Confirmed facts}




\begin{center}
\scriptsize
\begin{longtable}{>{\raggedright\arraybackslash}p{0.455\linewidth} >{\raggedright\arraybackslash}p{0.455\linewidth}}
\toprule
Field & Data \\
\midrule
\endfirsthead
\toprule
Field & Data \\
\midrule
\endhead
Source type & Portfolio case study \\
Company & Virtual Verse Studio \\
Region & Egypt / USA \\
Industry & Rail / training \\
Technology & VR training \\
Project type & Scenario-based training \\
Access model & VR \\
Metrics & Measurable behavior improvement; no numeric data \\
Confidence & Medium \\
\bottomrule
\end{longtable}
\end{center}




\subsubsection{Problem pattern}




Traditional role-play and lectures lacked realism for high-stress scenarios.




\subsubsection{Solution pattern}




The platform used first-person VR, perspective switching, behavioral design, and structured debriefs.




\subsubsection{Direct relevance to TwinSight}




The domain is weaker, but the \textbf{training logic} is useful. TwinSight can borrow the idea of scenario-driven learning and debrief-like explanation.




\subsubsection{What Alexander can adapt}




\begin{itemize}
\item Add a “User task flow” section.
\item Explain how the viewer supports comprehension, not only viewing.
\item Include user evaluation results.
\item Use Think-Aloud findings as qualitative validation.
\end{itemize}




\subsubsection{What not to copy}




\begin{itemize}
\item Do not claim behavioral training.
\item Do not claim soft-skill outcomes.
\item Do not make TwinSight sound like a rail-training product.
\end{itemize}




\medskip\hrule\medskip




\section{4. Cross-case patterns}




\subsection{4.1 Common case-study structure}




Across the strongest examples, the structure is:




\begin{mdcode}
Bloque de codigo: text
1. Title
2. One-line value proposition
3. Hero image or video
4. Problem / challenge
5. Solution
6. Key features
7. Results / impact
8. Technical snapshot
9. Delivery / process
10. Related work or CTA
\end{mdcode}




\subsection{4.2 Recommended TwinSight case-study structure}




\begin{mdcode}
Bloque de codigo: text
1. Title
2. One-sentence value proposition
3. Demo video / hero GIF
4. Problem: 2D documentation fails spatial reasoning
5. Solution: Unity WebGL technical visualization
6. Key features
7. Technical implementation
8. CAD-to-realtime pipeline
9. Optimization results
10. Evaluation results
11. Role and contribution
12. Constraints and limitations
13. Next steps
14. Links
\end{mdcode}




\subsection{4.3 Feature framing pattern}




Do not list features as isolated technical modules.




Frame them as user-facing workflows.




\begin{center}
\scriptsize
\begin{longtable}{>{\raggedright\arraybackslash}p{0.455\linewidth} >{\raggedright\arraybackslash}p{0.455\linewidth}}
\toprule
Technical feature & Better case-study framing \\
\midrule
\endfirsthead
\toprule
Technical feature & Better case-study framing \\
\midrule
\endhead
Component selection & Identify parts and relationships \\
Exploded view & Understand assembly hierarchy \\
Cross-section / clipping & Inspect hidden internal structures \\
Visual modes & Compare technical interpretations \\
Technical UI & Access part data while inspecting \\
CAD optimization & Make complex CAD browser-ready \\
WebGL deployment & Remove install friction \\
SUS / NASA-TLX & Validate usability and cognitive load \\
\bottomrule
\end{longtable}
\end{center}




\medskip\hrule\medskip




\section{5. Metrics pattern}




\subsection{5.1 Metrics seen in collected examples}




\begin{center}
\scriptsize
\begin{longtable}{>{\raggedright\arraybackslash}p{0.455\linewidth} >{\raggedright\arraybackslash}p{0.455\linewidth}}
\toprule
Metric type & Example \\
\midrule
\endfirsthead
\toprule
Metric type & Example \\
\midrule
\endhead
Engagement & 3× engagement \\
Lead quality & 45\% more qualified leads \\
Time reduction & 2-week shorter cycle \\
Downtime & 30\% reduction \\
Planning time & 50\% reduction \\
Comprehension & Improved engagement/comprehension \\
Behavior & Measurable behavior change \\
Decision speed & Faster decision cycles \\
\bottomrule
\end{longtable}
\end{center}




\subsection{5.2 Metrics available for TwinSight}




TwinSight has unusually strong validation data for a portfolio project:




\begin{center}
\scriptsize
\begin{longtable}{>{\raggedright\arraybackslash}p{0.455\linewidth} >{\raggedright\arraybackslash}p{0.455\linewidth}}
\toprule
Metric & TwinSight value \\
\midrule
\endfirsthead
\toprule
Metric & TwinSight value \\
\midrule
\endhead
Optimized triangles & 95,617 \\
Original CAD routes & 6.5M+ triangles \\
Participants & 12 \\
SUS average & 91.88 \\
NASA-TLX Raw, 3D viewer & 8.69 \\
NASA-TLX Raw, 2D support & 19.89 \\
Task-condition records & 96 \\
\bottomrule
\end{longtable}
\end{center}




\subsection{5.3 How to present metrics}




Use direct, defensible phrasing:




\begin{mdcode}
Bloque de codigo: text
Optimized CAD-derived drone geometry from multi-million-triangle source routes into a 95,617-triangle WebGL-ready model.
\end{mdcode}




\begin{mdcode}
Bloque de codigo: text
Evaluated the viewer with 12 participants using SUS, NASA-TLX Raw, and Think-Aloud protocols.
\end{mdcode}




\begin{mdcode}
Bloque de codigo: text
Observed lower perceived workload in the 3D viewer condition compared with 2D support.
\end{mdcode}




Do not overstate causality.




Avoid:




\begin{mdcode}
Bloque de codigo: text
Proved that 3D is better than 2D.
\end{mdcode}




Use:




\begin{mdcode}
Bloque de codigo: text
The formative evaluation suggested lower perceived workload and high usability for the interactive 3D viewer.
\end{mdcode}




\medskip\hrule\medskip




\section{6. Visual pattern benchmark}




\subsection{6.1 Media types used}




\begin{center}
\scriptsize
\begin{longtable}{>{\raggedright\arraybackslash}p{0.305\linewidth} >{\raggedright\arraybackslash}p{0.305\linewidth} >{\raggedright\arraybackslash}p{0.305\linewidth}}
\toprule
Media type & Seen in examples & Recommended for TwinSight \\
\midrule
\endfirsthead
\toprule
Media type & Seen in examples & Recommended for TwinSight \\
\midrule
\endhead
Hero image & Yes & Yes \\
Demo video & Yes & Yes \\
Screenshot sequence & Yes & Yes \\
Feature cards & Yes & Yes \\
Process diagram & Sometimes & Yes \\
Before/after & Often implied & Yes \\
Metrics cards & Yes & Yes \\
Client quote & Sometimes & Replace with thesis/evaluation quote \\
CTA & Yes & Yes \\
\bottomrule
\end{longtable}
\end{center}




\subsection{6.2 Required TwinSight visuals}




\begin{center}
\scriptsize
\begin{longtable}{>{\raggedright\arraybackslash}p{0.455\linewidth} >{\raggedright\arraybackslash}p{0.455\linewidth}}
\toprule
Visual & Purpose \\
\midrule
\endfirsthead
\toprule
Visual & Purpose \\
\midrule
\endhead
Hero drone viewer shot & Immediate comprehension \\
Exploded view shot & Assembly understanding \\
Clipping/cross-section shot & Technical inspection \\
Component selection shot & Interaction proof \\
Visual modes grid & Technical art proof \\
CAD/high-poly/optimized comparison & Pipeline proof \\
Triangle reduction card & Optimization proof \\
SUS/NASA-TLX card & Validation proof \\
Mobile/web view & WebGL deployment proof \\
Architecture/pipeline diagram & Developer credibility \\
\bottomrule
\end{longtable}
\end{center}




\medskip\hrule\medskip




\section{7. Language patterns to adopt}




\subsection{Strong phrases from the benchmark}




Use adapted versions of these patterns:




\begin{mdcode}
Bloque de codigo: text
browser-based 3D
real-time customization
interactive technical visualization
spatial context
navigable 3D replica
technical inspection
structured user flow
web-accessible experience
no-install deployment
measurable validation
faster understanding
component-level exploration
\end{mdcode}




\subsection{Phrases to avoid}




Avoid unless directly true:




\begin{mdcode}
Bloque de codigo: text
live IoT integration
predictive maintenance
enterprise deployment
production-ready digital twin
WebAR placement
CRM integration
safety outcome
behavior change
client revenue impact
industrial metaverse
\end{mdcode}




\medskip\hrule\medskip




\section{8. Recommended TwinSight case-study copy skeleton}




\subsection{Title}




\begin{mdcode}
Bloque de codigo: text
TwinSight X500 — Unity WebGL Technical Visualization for Drone Assembly Inspection
\end{mdcode}




\subsection{Subtitle}




\begin{mdcode}
Bloque de codigo: text
An interactive browser-based 3D viewer that helps users understand the spatial relationships between drone components through selection, exploded view, clipping, visual modes, and technical part panels.
\end{mdcode}




\subsection{Problem}




\begin{mdcode}
Bloque de codigo: text
Traditional 2D assembly documentation forces users to mentally reconstruct spatial relationships between parts. For complex drone assemblies, this increases cognitive effort and makes hidden structures, orientation, and component hierarchy harder to understand.
\end{mdcode}




\subsection{Solution}




\begin{mdcode}
Bloque de codigo: text
TwinSight X500 transforms CAD-derived drone geometry into an optimized Unity WebGL experience. Users can inspect components, isolate assemblies, activate exploded views, cut through structures, switch visual modes, and access technical information directly in the browser.
\end{mdcode}




\subsection{Key features}




\begin{mdcode}
Bloque de codigo: text
- Component selection and highlighting
- Exploded assembly view
- Cross-section / clipping inspection
- Technical part information panels
- Realistic, X-Ray, blueprint, ghosted, wireframe, solid, and thermal-style modes
- CAD-to-realtime geometry optimization
- Mobile-first WebGL interaction design
- Usability and workload evaluation
\end{mdcode}




\subsection{Results}




\begin{mdcode}
Bloque de codigo: text
- Reduced CAD-derived geometry from routes exceeding 6.5M triangles to 95,617 optimized triangles
- Evaluated with 12 participants
- SUS average: 91.88
- NASA-TLX Raw: 8.69 for 3D viewer vs 19.89 for 2D support
- 96 task-condition records
\end{mdcode}




\subsection{Role}




\begin{mdcode}
Bloque de codigo: text
I developed the Unity WebGL integration, runtime interaction systems, technical UI, visualization modes, asset optimization workflow, Blender-based preparation, documentation, and evaluation pipeline, using AI assistance as an implementation support tool while retaining responsibility for architecture, debugging, integration, validation, and final technical decisions.
\end{mdcode}




\medskip\hrule\medskip




\section{9. TwinSight case-study section order}




\subsection{Recommended final order}




\begin{mdcode}
Bloque de codigo: text
01. Hero
02. One-line value proposition
03. Problem
04. Solution
05. Demo video
06. Key features
07. Technical pipeline
08. CAD-to-realtime optimization
09. Interaction systems
10. Visual modes
11. Evaluation results
12. Contribution and AI-assisted workflow
13. Constraints and limitations
14. Next steps
15. Links
\end{mdcode}




\subsection{Why this order}




The benchmark case studies show that readers need value first, then details.




Do not start with academic methodology.




Start with:




\begin{mdcode}
Bloque de codigo: text
What problem did this solve?
What can the user do?
What technical proof exists?
What evidence validates it?
\end{mdcode}




Then include academic details deeper.




\medskip\hrule\medskip




\section{10. Gaps in current benchmark}




\subsection{Missing categories}




The current dataset lacks enough examples from:




\begin{center}
\scriptsize
\begin{longtable}{>{\raggedright\arraybackslash}p{0.455\linewidth} >{\raggedright\arraybackslash}p{0.455\linewidth}}
\toprule
Category & Status \\
\midrule
\endfirsthead
\toprule
Category & Status \\
\midrule
\endhead
Drone visualization & Missing \\
Robotics visualization & Missing \\
CAD-to-Unity case studies & Weak \\
Unity WebGL technical inspection tools & Weak \\
Medical 3D visualization & Missing \\
Aerospace assembly visualization & Weak \\
Three.js technical viewers & Missing \\
WebGPU industrial viewers & Missing \\
Academic HCI visualization demos & Missing \\
\bottomrule
\end{longtable}
\end{center}




\subsection{Suggested expansion query set}




Use these later:




\begin{mdcode}
Bloque de codigo: text
"CAD to Unity case study"
"Unity technical visualization case study"
"WebGL 3D inspection tool case study"
"drone 3D visualization case study"
"robotics digital twin Unity case study"
"Three.js product configurator case study"
"interactive assembly instructions 3D case study"
"WebGL technical viewer CAD"
"Unity digital twin manufacturing case study"
"3D maintenance training Unity case study"
\end{mdcode}




\medskip\hrule\medskip




\section{11. How this changes the roadmap}




\subsection{No major strategy change}




The case-study benchmark confirms the existing direction:




\begin{mdcode}
Bloque de codigo: text
Unity WebGL + technical visualization + CAD-to-realtime + browser-based interaction
\end{mdcode}




remains the strongest route.




\subsection{Stronger emphasis needed}




TwinSight should be presented less like:




\begin{mdcode}
Bloque de codigo: text
academic thesis viewer
\end{mdcode}




and more like:




\begin{mdcode}
Bloque de codigo: text
technical visualization case study with measured usability and workload evaluation
\end{mdcode}




\subsection{Best market language}




Use:




\begin{mdcode}
Bloque de codigo: text
technical visualization
interactive 3D
browser-based 3D
CAD-to-realtime
component inspection
spatial understanding
digital-twin-adjacent
training and assembly support
\end{mdcode}




Use cautiously:




\begin{mdcode}
Bloque de codigo: text
digital twin
simulation
XR
technical training
\end{mdcode}




Avoid as primary:




\begin{mdcode}
Bloque de codigo: text
game project
3D art project
academic demo
prototype only
AI-built app
\end{mdcode}




\medskip\hrule\medskip




\section{12. Required updates to existing files}




\subsection{Update \texttt{08\_twinsight\_x500\_case\_study.md}}




Add:




\begin{itemize}
\item problem/solution/results structure;
\item short value proposition;
\item feature cards;
\item metrics cards;
\item source-inspired case-study layout;
\item explicit “not a full digital twin” limitation;
\item digital-twin-adjacent language.
\end{itemize}




\subsection{Update \texttt{20\_portfolio\_copy\_and\_site\_structure.md}}




Add:




\begin{itemize}
\item case-study page structure;
\item portfolio card copy;
\item feature labels;
\item CTA labels;
\item metrics section.
\end{itemize}




\subsection{Update \texttt{21\_twinsight\_demo\_video\_script\_and\_shotlist.md}}




Add:




\begin{itemize}
\item first 5 seconds: value proposition;
\item feature sequence based on problem-solution-result;
\item metrics card near end;
\item final links.
\end{itemize}




\subsection{Update \texttt{17\_cv\_base\_and\_role\_variants.md}}




Add one case-study bullet pattern:




\begin{mdcode}
Bloque de codigo: text
Built a Unity WebGL technical visualization prototype for drone assembly inspection, transforming CAD-derived assets into an optimized browser-based viewer with component selection, exploded view, clipping, visual modes, and usability/workload evaluation.
\end{mdcode}




\medskip\hrule\medskip




\section{13. Next task recommendation}




The next task should be:




\begin{mdcode}
Bloque de codigo: text
F1_demo_reels_benchmark
\end{mdcode}




\subsection{Tool recommendation}




\begin{center}
\scriptsize
\begin{longtable}{>{\raggedright\arraybackslash}p{0.455\linewidth} >{\raggedright\arraybackslash}p{0.455\linewidth}}
\toprule
Tool & Use \\
\midrule
\endfirsthead
\toprule
Tool & Use \\
\midrule
\endhead
Agent mode & Yes \\
Deep Research & No \\
\bottomrule
\end{longtable}
\end{center}




Reason:




The current file defines what the case study should say. The next missing piece is how the \textbf{demo video} should be edited and sequenced.




\subsection{Prompt for next task}




\begin{mdcode}
Bloque de codigo: text
Execute F1_demo_reels_benchmark.

Objective:
Collect public demo reel and project demo examples relevant to TwinSight X500:
technical artist demo reels, Unity technical reels, real-time 3D developer reels, technical visualization demos, XR/simulation demos, CAD-to-realtime/product visualization demos, WebGL interactive 3D demos.

Find 15–25 examples.

Output structured tables only.

Fields:
- Demo label
- URL
- Platform
- Role target
- Duration
- Opening structure
- Shot sequence
- Captions/text usage
- Voiceover yes/no
- Technical breakdown yes/no
- Tools shown
- Metrics shown
- Call-to-action
- What Alexander can adapt
- What Alexander should not copy
- Confidence
- Date checked

Rules:
- Use only public videos/pages.
- Do not invent duration if not visible.
- If unavailable, write “not found”.
- Keep cells short.
- Tables only.
- Include URLs.
- No strategy.
\end{mdcode}



